\uuid{4nT7}
\titre{Série télescopique cachée et équivalent }

\niveau{L1}
\module{Séries numériques}
\chapitre{Séries à termes positifs}
\sousChapitre{Séries télescopiques et équivalents}
\theme{Fonction arctangente et nature d'une série}
\auteur{Quentin Liard}
\datecreate{ 2026-04-10}
\organisation{AMSCC}
\difficulte{3}

\contenu{

\texte{ 
Pour tout entier $n\geq 1$, on pose
$$
u_n=\arctan\left(\frac1n\right)-\arctan\left(\frac{1}{n+1}\right).
$$
On considère la série
$
\sum_{n\geq 1} u_n,
$
et on rappelle que pour tout $x$ réel,
 $$\arctan'(x)=\frac{1}{1+x^2}.$$
}

\begin{enumerate}

\item   \question{Montrer que, pour tout $n\geq 1$, on a $u_n\geq 0$.}
\reponse{
La fonction $\arctan$ est strictement croissante sur $\mathbb{R}$.
Or, pour tout $n\geq 1$,
$
\frac1n>\frac{1}{n+1}.
$
Donc
$
\arctan\left(\frac1n\right)>\arctan\left(\frac{1}{n+1}\right),
$
et par conséquent
$
u_n\geq 0.
$
}

\item  \question{Donner le développement en série entière en $0$ de $x \mapsto \frac{1}{1+x^2}$.}
\reponse{
On part de la série géométrique classique
$$
\frac{1}{1-t}=\sum_{n=0}^{+\infty} t^n,
\qquad |t|<1.
$$
En prenant $t=-x^2$, on obtient
$$
\frac{1}{1+x^2}=\sum_{n=0}^{+\infty}(-1)^n x^{2n},
\qquad |x|<1.
$$
Ainsi, le développement en série entière en $0$ est
$$
\boxed{\frac{1}{1+x^2}=\sum_{n=0}^{+\infty}(-1)^n x^{2n}}
\qquad \text{pour } |x|<1.
$$
}

\item \question{En déduire le développement en série entière en $0$ de $x\mapsto \arctan(x)$.}
\reponse{
Comme
$$
\arctan'(x)=\frac{1}{1+x^2}
$$
et que
$$
\frac{1}{1+x^2}=\sum_{n=0}^{+\infty}(-1)^n x^{2n}
\qquad \text{pour } |x|<1,
$$
on peut intégrer terme à terme entre $0$ et $x$ :
$$
\arctan(x)-\arctan(0)=\int_0^x \frac{1}{1+t^2}\,dt
=\int_0^x \sum_{n=0}^{+\infty}(-1)^n t^{2n}\,dt.
$$
Donc
$$
\arctan(x)=\sum_{n=0}^{+\infty}(-1)^n \frac{x^{2n+1}}{2n+1},
\qquad |x|<1.
$$
Ainsi,
$$
\boxed{\arctan(x)=\sum_{n=0}^{+\infty}(-1)^n \frac{x^{2n+1}}{2n+1}}
\qquad \text{pour } |x|<1.
$$
}

\item \question{Démontrer que pour $x$ proche de $0$, $\arctan(x)=x+o(x^2)$.}
\reponse{
D’après la question précédente,
$$
\arctan(x)=x-\frac{x^3}{3}+\frac{x^5}{5}-\cdots
$$
au voisinage de $0$. On en déduit
$$
\arctan(x)-x=-\frac{x^3}{3}+o(x^3).
$$
En particulier,
$$
\arctan(x)-x=o(x^2).
$$
Donc
$$
\boxed{\arctan(x)=x+o(x^2)}
\qquad (x\to 0).
$$
}

\item   \question{Déterminer un équivalent simple de $u_n$ lorsque $n\to+\infty$.}
\reponse{
En utilisant le résultat précédent,
$$
\arctan\left(\frac1n\right)=\frac1n+o\left(\frac1{n^2}\right)
\qquad\text{et}\qquad
\arctan\left(\frac1{n+1}\right)=\frac1{n+1}+o\left(\frac1{n^2}\right).
$$
Donc
$$
u_n
=
\left(\frac1n-\frac1{n+1}\right)+o\left(\frac1{n^2}\right).
$$
Or
$$
\frac1n-\frac1{n+1}=\frac{1}{n(n+1)}\sim \frac1{n^2}.
$$
Ainsi,
$$
\boxed{u_n\sim \frac1{n^2}}.
$$
}

\item   \question{Établir la convergence de la série $\sum u_n$ en la comparant à une série de Riemann.}
\reponse{
D’après la question précédente,
$
u_n\sim \frac{1}{n^2}.
$
La série
$
\sum_{n\geq 1}\frac{1}{n^2}
$
est une série de Riemann convergente. Donc, par comparaison par équivalence, la série
$
\sum_{n\geq 1}u_n
$
converge.
}

\item   \question{Montrer que la série $\sum u_n$ est télescopique et calculer sa somme.}
\reponse{
La somme partielle d’ordre $N$ vaut
$$
S_N=\sum_{n=1}^N \left(\arctan\left(\frac1n\right)-\arctan\left(\frac{1}{n+1}\right)\right).
$$
Les termes se simplifient deux à deux :
$
S_N=\arctan(1)-\arctan\left(\frac{1}{N+1}\right).
$
Comme
$
\arctan(1)=\frac{\pi}{4}
\qquad\text{et}\qquad
\arctan\left(\frac{1}{N+1}\right)\xrightarrow[N\to+\infty]{}0,
$
on obtient
$
\sum_{n\geq 1}u_n=\frac{\pi}{4}.
$
La série converge donc, et sa somme vaut
$
\boxed{\frac{\pi}{4}}.
$
}
\end{enumerate}

}
